\chapter{Architecture technique}

\section{Langage et interface}

L'application est écrite exclusivement en Kotlin, avec une interface
entièrement en Jetpack Compose : aucun layout XML. Elle repose sur une seule
activité et Navigation Compose. Le thème Material~3 sert de base, avec les
couleurs et les polices de la direction C.

\begin{tableau}{L{4cm} Y}
  \entete{Paramètre} & \entete{Valeur} \\
  \midrule
  Version minimale & Android 8.0 (API 26) \\
  \midrule
  Version cible & Dernière API stable au moment du rendu \\
  \midrule
  Orientation & Portrait ; paysage pour le mode cuisine \\
\end{tableau}

\section{Architecture MVVM}

Chaque écran suit le même schéma. Le composable affiche un état et remonte
les actions de l'utilisateur. Le ViewModel transforme ces actions en appels
au repository et expose un \texttt{StateFlow} d'état d'écran. Le
repository est la seule porte d'accès aux données : il lit la base locale et
la rafraîchit depuis le réseau.

\begin{figure}[htbp]
  \centering
  \begin{tikzpicture}[
      couche/.style = {rounded corners=3mm, fill=surface, text=texte,
                       minimum width=3.6cm, minimum height=1.2cm,
                       align=center, font=\small\policeMoyenne},
      source/.style = {rounded corners=3mm, draw=filet, line width=0.6pt,
                       text=texte, minimum width=3.3cm, minimum height=1cm,
                       align=center, font=\small},
      fleche/.style = {-{Latex[length=2mm]}, draw=texte2, line width=0.6pt},
      note/.style   = {font=\scriptsize, text=texte2},
    ]
    \node[couche] (ui) {Écran\\{\footnotesize Composable}};
    \node[couche, right=1.6cm of ui] (vm) {ViewModel};
    \node[couche, right=1.6cm of vm, fill=teinteAction] (repo) {Repository};

    \draw[fleche] ([yshift=3mm]ui.east) -- node[note, above]{actions} ([yshift=3mm]vm.west);
    \draw[fleche] ([yshift=-3mm]vm.west) -- node[note, below]{état} ([yshift=-3mm]ui.east);
    \draw[fleche] ([yshift=3mm]vm.east) -- node[note, above]{appels} ([yshift=3mm]repo.west);
    \draw[fleche] ([yshift=-3mm]repo.west) -- node[note, below]{Flow} ([yshift=-3mm]vm.east);

    \node[source, below=1.6cm of ui]   (api)   {API REST\\{\footnotesize Retrofit}};
    \node[source, right=0.55cm of api] (room)  {Room\\{\footnotesize cache, favoris}};
    \node[source, right=0.55cm of room](store) {DataStore\\{\footnotesize préférences}};
    \node[source, right=0.55cm of store](json) {Assets JSON\\{\footnotesize vins, règles}};

    \foreach \s in {api, room, store, json}
      \draw[fleche] (repo.south) -- (\s.north);
  \end{tikzpicture}
  \caption{Flux de données entre les couches}
  \label{fig:mvvm}
\end{figure}

\begin{tableau}{L{4cm} Y}
  \entete{Besoin} & \entete{Bibliothèque} \\
  \midrule
  Injection de dépendances & Hilt \\
  \midrule
  Réseau & Retrofit, OkHttp, kotlinx.serialization \\
  \midrule
  Images & Coil, avec cache disque \\
  \midrule
  Navigation & Navigation Compose \\
  \midrule
  Base locale & Room \\
  \midrule
  Préférences & DataStore \\
  \midrule
  Asynchronisme & Coroutines et Flow \\
\end{tableau}

Le code tient dans un seul module Gradle, découpé en paquets par
fonctionnalité (\texttt{onboarding}, \texttt{carte}, \texttt{explorer},
\texttt{recette}, \texttt{accords}, \texttt{courses}, \texttt{favoris},
\texttt{reglages}) et en paquets communs (\texttt{designsystem},
\texttt{data}). À deux, plusieurs modules apporteraient plus de
configuration que de bénéfices.

\section{Sources de données}

\subsection{API}

Les deux API sont publiques et gratuites. La clé de test « 1 » suffit pour un
projet pédagogique.

\begin{tableau}{L{2.6cm} Y}
  \entete{API} & \entete{Points d'accès utilisés} \\
  \midrule
  TheMealDB &
    \texttt{\small search.php?s=} (nom),
    \texttt{\small filter.php?c= / a= / i=} (catégorie, pays, ingrédient),
    \texttt{\small lookup.php?i=} (fiche),
    \texttt{\small random.php},
    \texttt{\small categories.php},
    \texttt{\small list.php?a=list} (pays) \\
  \midrule
  TheCocktailDB &
    \texttt{\small search.php?s=},
    \texttt{\small filter.php?a=} (avec ou sans alcool),
    \texttt{\small filter.php?g=} (verre),
    \texttt{\small lookup.php?i=},
    \texttt{\small random.php} \\
\end{tableau}

\subsection{Données ajoutées par l'équipe}

Une partie de ce qu'affichent les maquettes n'existe dans aucune des deux
API. Ces données sont soit rédigées par l'équipe et embarquées en JSON dans
l'application, soit déduites du contenu des recettes.

\begin{tableau}{L{3.6cm} Y}
  \entete{Donnée} & \entete{Origine} \\
  \midrule
  Vins et cépages & Catalogue JSON rédigé par l'équipe \\
  \midrule
  Règles d'accord & Table JSON : catégorie et pays du plat vers styles de vin et de cocktail \\
  \midrule
  Allergènes et substitutions & Table JSON des ingrédients courants (gluten, lactose, produits animaux) \\
  \midrule
  Rayons de la liste & Table JSON ingrédient vers rayon \\
  \midrule
  Régions et plat emblématique & Table JSON pour la carte du monde \\
  \midrule
  Matériel et durées & Mots-clés repérés dans les instructions (four, minutes…) \\
  \midrule
  Portions & Quatre personnes par défaut, quantités numériques extraites des mesures \\
\end{tableau}

\section{Stockage local et cache}

Room est la source de vérité : les écrans observent la base et le réseau ne
fait que la mettre à jour.

\begin{tableau}{L{3.6cm} Y}
  \entete{Table} & \entete{Contenu} \\
  \midrule
  \texttt{recette}, \texttt{boisson} & Fiches consultées, avec date de mise en cache \\
  \midrule
  \texttt{carte\_du\_jour} & Date, plat, verre et dessert tirés \\
  \midrule
  \texttt{favori} & Type (plat, boisson, menu), identifiant, date \\
  \midrule
  \texttt{menu}, \texttt{menu\_element} & Menus composés et leurs éléments \\
  \midrule
  \texttt{article\_courses} & Nom, quantité, unité, rayon, état coché \\
\end{tableau}

Les préférences (âge confirmé, régime, matériel, boissons, thème, retour
haptique, profil) sont stockées dans DataStore.

Les fiches ouvertes sont conservées 30 jours, et sans limite tant qu'elles
figurent dans un favori, un menu ou la carte du jour. Les résultats de
recherche ne sont pas mis en cache. Les photos passent par le cache disque
de Coil.

\section{Asynchronisme et erreurs réseau}

Les appels réseau et les accès à la base s'exécutent dans des coroutines,
hors du thread principal. Chaque écran expose trois états : chargement,
contenu, erreur. Les erreurs sont traduites en messages compréhensibles :

\begin{tableau}{L{4cm} Y}
  \entete{Situation} & \entete{Comportement} \\
  \midrule
  Pas de réseau & Bandeau « hors ligne », affichage des données en cache \\
  \midrule
  Délai dépassé (15 s) & Message et bouton Réessayer \\
  \midrule
  Erreur du serveur & Message et bouton Réessayer, données en cache conservées \\
  \midrule
  Réponse vide & État vide avec une suggestion (élargir les filtres) \\
\end{tableau}
